\chapter{随机过程与随机逼近}
\label{chap:stochastic-processes-and-approximation}

第~\ref{chap:probability-theory}章考察独立观测的平均。若每次观测都来自同一个分布，
大数定律说明样本均值是否接近期望，集中不等式进一步控制有限样本误差。然而，许多学习
过程会把当前结果带到下一步：采样器从当前位置产生新状态，算法根据已有数据选择下一次
查询，参数又用刚得到的噪声观测继续更新。此时，记录数仍在增加，独立性却已经消失。

本章围绕一个均值求根问题研究这种时间结构。设长期观测均值为$\mu$，目标是求方程
\begin{equation}
  h(\theta)=\mathbb{E}[Y-\theta]=\mu-\theta=0
  \label{eq:stoch-mean-root}
\end{equation}
的根$\theta^\star=\mu$。若观测$Y_t$来自Markov链，它们可能彼此相关，因而不能预设
独立性；若下一次观测方向由过去结果决定，误差还会与设计共同演化；若不保存全部数据而
直接更新$\theta^{(t)}$，问题又从估计固定平均变成带噪迭代。本章以均值求根作为贯穿
示例，但各节并不估计同一个参数：具体目标会从长期均值扩展到线性参数和一般平均漂移的
根。四种视角共享的是如何在时间依赖或自适应信息下分离平均信号与噪声，其中鞅提供连接
这些任务的条件噪声工具。

我们先判断相关序列的时间平均何时仍指向$\mu$，再用鞅刻画给定过去后无法预测的噪声。
条件方差给出比统一增量界更细的序贯控制，带可预测特征的向量鞅则把误差放进随数据增长的
设计矩阵。最后，随机逼近把这些噪声工具接到求根更新，并说明步长、平均漂移与迭代有界性
各自承担什么作用。贯穿各节的重点不是把“独立样本数”换成“时间步数”，而是辨认过去信息
怎样改变当前分布，以及证明中哪一步仍能利用条件平均为零。

\section{Markov链与长期平均}
\label{sec:stoch-markov-long-run-average}

\subsection{转移规则、平稳分布与时间齐次性}

设$(Z_t)_{t\geq0}$取值于可测空间$(\mathcal{Z},\mathcal{A})$。若对每个
$B\in\mathcal{A}$，
\[
  \mathbb{P}(Z_{t+1}\in B\mid Z_0,\ldots,Z_t)
  =
  \mathbb{P}(Z_{t+1}\in B\mid Z_t),
\]
则称$(Z_t)$为\term{Markov链}（Markov chain）。这个条件不是说过去没有影响，而是
说当前状态已经保存了预测下一步所需的全部历史信息。若右侧可统一写成
$P(Z_t,B)$，且转移核$P$不随$t$改变，则链是\term{时间齐次的}
（time-homogeneous）。时间齐次描述机制不变，不代表$Z_t$的边缘分布不随时间改变
\scite{math-meyn2009markov}

这里的转移核需要同时满足两项可测性条件：固定$z$时，$P(z,\cdot)$是
$(\mathcal{Z},\mathcal{A})$上的概率测度；固定$B\in\mathcal{A}$时，
$P(\cdot,B)$是可测函数。对有界可测函数$f$，记
\[
  Pf(z)=\int_{\mathcal{Z}}f(z')P(z,\mathrm{d}z').
\]
两个核的复合按
\[
  (PQ)(z,B)
  =
  \int_{\mathcal{Z}}Q(z',B)P(z,\mathrm{d}z')
\]
定义；它表示先按$P$走一步，再按$Q$走一步。令$P^0(z,B)=\mathbb{I}\{z\in B\}$，
$P^{m+1}=P^mP$，则时间齐次链满足
\begin{equation}
  \mathbb{P}(Z_{t+m}\in B\mid Z_t=z)=P^m(z,B),
  \qquad
  P^{m+n}(z,B)
  =
  \int_{\mathcal{Z}}P^n(z',B)P^m(z,\mathrm{d}z').
  \label{eq:stoch-chapman-kolmogorov}
\end{equation}
第二个等式称为\term{Chapman--Kolmogorov关系}。证明可直接使用条件概率的塔式性质：
给定$Z_{t+m}=z'$后，Markov性质把后$n$步的条件概率化为$P^n(z',B)$；再对
$Z_{t+m}$在给定$Z_t=z$时的分布$P^m(z,\mathrm{d}z')$积分，就得到右侧。取$n=1$
并归纳，还同时证明了第一个等式及核幂确实给出多步转移。

概率分布$\pi$若满足
\begin{equation}
  \pi(B)=\int_{\mathcal{Z}}P(z,B)\,\pi(\mathrm{d}z),
  \qquad B\in\mathcal{A},
  \label{eq:stoch-stationary-distribution}
\end{equation}
就称为$P$的\term{平稳分布}（stationary distribution）。若$Z_0\sim\pi$，式
\eqref{eq:stoch-stationary-distribution}保证每个$Z_t$仍服从$\pi$；若从指定状态
$z_0$启动，平稳性本身却没有说明$\mathcal{L}(Z_t)$会趋向$\pi$。恒等转移核
$P(z,B)=\mathbb{I}\{z\in B\}$保持任意分布，但它永远不忘记初态，这已经表明“保持
目标”和“到达目标”是两项不同要求。

“平稳分布”是核与一个概率分布之间的关系，\term{平稳过程}（stationary process）则要求
任意有限维联合分布在时间平移后不变。对时间齐次Markov链，若$Z_0\sim\pi$且
$\pi P=\pi$，则
\[
  \mathbb{P}(Z_0\in\mathrm{d}z_0,\ldots,Z_m\in\mathrm{d}z_m)
  =
  \pi(\mathrm{d}z_0)
  \prod_{j=1}^{m}P(z_{j-1},\mathrm{d}z_j).
\]
把起点向后平移一格时，先对$z_0$积分，$\pi P=\pi$把$Z_1$的边缘重新化为$\pi$，
其余转移核不变；反复应用便证明所有有限维分布都不随平移改变。因此，从平稳分布启动的
时间齐次链是平稳过程。只知道某一个时刻的边缘分布等于$\pi$，却没有核不变或联合分布
条件时，不能据此断言整个过程平稳。

一个常用的保持性条件是\term{细致平衡}（detailed balance）
\begin{equation}
  \pi(\mathrm{d}z)P(z,\mathrm{d}z')
  =\pi(\mathrm{d}z')P(z',\mathrm{d}z).
  \label{eq:stoch-detailed-balance}
\end{equation}
它要求平稳运行时任意两处之间的概率流量正反相等。对$z$积分便得到
\eqref{eq:stoch-stationary-distribution}，所以细致平衡足以推出平稳性。反向不成立：
围绕三个状态单向循环的链可以保持均匀分布，却没有逐边对称的反向流。细致平衡是构造
采样链的便利证书，不是遍历平均成立的必要定义。

\begin{example}[相关的两状态观测]
  \label{ex:stoch-two-state-chain}
  令$Z_t\in\{0,2\}$，并把观测直接取为$Y_t=Z_t$。转移矩阵按状态$0,2$排序为
  \[
    \bm{P}
    =
    \begin{pmatrix}
      0.9 & 0.1\\
      0.2 & 0.8
    \end{pmatrix}.
  \]
  平稳分布为$\bm{\pi}=(2/3,1/3)^\top$，因为
  $\bm{P}^{\top}\bm{\pi}=\bm{\pi}$。长期均值因此是
  \[
    \mu=\mathbb{E}_{\pi}[Y]=\frac23.
  \]
  两条跨状态概率流都等于
  $(2/3)\times0.1=(1/3)\times0.2=1/15$，故该链还满足细致平衡。
  从状态$0$启动时，$Y_0$不服从$\pi$；转移规则时间齐次，过程却要经过一段时间才
  接近平稳。
\end{example}

\subsection{可达性、周期与收敛不是同一条件}

在有限状态空间中，\term{不可约性}（irreducibility）要求任意状态都能以正概率在有限步
内到达任意另一状态。它排除多个互不连通的封闭类，并保证平稳分布唯一。
\term{周期}（period）是从某状态返回自身的可能步数集合的最大公约数；不可约链的所有
状态具有相同周期。周期等于$1$时，链称为\term{非周期的}（aperiodic）。

更具体地，状态$i$到$j$可达，是指存在$m\geq0$使$P^m(i,\{j\})>0$。令
$T_i^+=\inf\{t\geq1:Z_t=i\}$；若
$\mathbb{P}_i(T_i^+<\infty)=1$，就称$i$为\term{常返状态}（recurrent state）。
有限链必有至少一个封闭常返类；不可约性迫使全部状态属于同一个类，所以每个状态都常返。
从$i$出发，把从一次访问$i$起、到下一次返回$i$前为止的路径作为一个游程；返回时刻处的
强Markov性说明这些游程独立同分布。记第$k$个游程的长度为$L_k$、其中访问$j$的次数为
$R_k$，则大数定律给出
\[
  \frac{\sum_{k=1}^{m}R_k}{\sum_{k=1}^{m}L_k}
  \longrightarrow
  \frac{\mathbb{E}_iR_1}{\mathbb{E}_iL_1}
  =\pi_j
  \quad\text{几乎必然成立。}
\]
长期时间占比来自多个游程的累计比值，而不是单次随机比值$R_k/L_k$。有限不可约性使
返回时间均值有限；取$j=i$时，每个游程恰计一次对$i$的访问，因而得到
$\pi_i=1/\mathbb{E}_iT_i^+$。这个返回分解是遍历平均的概率证明入口。矩阵方面，
$\bm{P}^{\top}$的特征值$1$对应唯一的概率右特征向量$\bm{\pi}$；等价地，
$\bm{\pi}^{\top}$是行随机转移矩阵$\bm{P}$的唯一概率左特征向量。加入非周期性后，
其他模态不会在单位圆上持续振荡，矩阵幂$\bm{P}^t$才逐行趋于
$\bm{\pi}^{\top}$。这两条路线分别解释了为什么周期不妨碍时间占比，却会妨碍逐时刻分布
收敛。

有限、不可约、非周期链满足
\begin{equation}
  P^t(z,\cdot)\longrightarrow\pi(\cdot)
  \quad\text{对每个初始状态 }z,
  \label{eq:stoch-finite-chain-distribution-convergence}
\end{equation}
其中有限状态下的收敛可以用\term{总变差距离}（total variation distance）衡量：
对概率测度$\nu,\pi$，定义
\[
  \lVert\nu-\pi\rVert_{\mathrm{TV}}
  =
  \sup_{B\text{ 可测}}|\nu(B)-\pi(B)|.
\]
不可约性与非周期性在这里服务于逐时刻分布收敛。时间平均的要求略有不同：有限不可约链
即使有周期，对每个函数
$f:\mathcal{Z}\to\mathbb{R}$仍有
\begin{equation}
  \frac1n\sum_{t=1}^{n}f(Z_t)\longrightarrow\mathbb{E}_{\pi}[f(Z)]
  \quad\text{几乎必然成立。}
  \label{eq:stoch-markov-ergodic-average}
\end{equation}
例如在$0$与$2$之间确定性交替的链具有周期$2$。从$0$出发时，$Z_t$的分布在两个点之间
来回切换，不存在逐步极限；但时间平均仍趋向$1$。因此，讨论“遍历”时必须说明结论针对
逐步分布、时间平均还是其他对象。

例~\ref{ex:stoch-two-state-chain}没有这些退化。其第二个特征值为$0.7$。令
$\widetilde{Y}_t=Y_t-\mu$，直接计算可得
\begin{equation}
  \mathbb{E}[\widetilde{Y}_{t+1}\mid Z_t]=0.7\widetilde{Y}_t.
  \label{eq:stoch-two-state-eigenfunction}
\end{equation}
反复使用塔式性质便得到
$\mathbb{E}[\widetilde{Y}_{t+k}\mid Z_t]=0.7^k\widetilde{Y}_t$。
若链从平稳分布启动，
\begin{equation}
  \operatorname{Corr}_{\pi}(Y_t,Y_{t+k})=0.7^k.
  \label{eq:stoch-two-state-autocorrelation}
\end{equation}
这个几何衰减同时展示了“收敛到平稳”和“样本之间仍相关”：边缘分布可以从一开始就等于
$\pi$，相邻观测却不会因此独立。

\subsection{混合速度、自相关与有效样本量}

\term{混合速度}（mixing rate）描述链忘记初态的快慢。有限链常考察
\[
  d(t)
  =
  \sup_{z\in\mathcal{Z}}
  \lVert P^t(z,\cdot)-\pi\rVert_{\mathrm{TV}},
\]
并把使$d(t)\leq\varepsilon$的最小$t$称为相应精度下的混合时间。这个最坏初态指标
针对分布距离；估计某个特定$f$的时间平均时，还要考察$f(Z_t)$自身的自相关。

若链平稳，记$\sigma_f^2=\operatorname{Var}_{\pi}(f(Z_0))$，
$\rho_f(k)=\operatorname{Corr}_{\pi}(f(Z_0),f(Z_k))$。展开样本均值方差得到
\begin{align}
  \operatorname{Var}\left(
    \frac1n\sum_{t=1}^{n}f(Z_t)
  \right)
  &=
  \frac{\sigma_f^2}{n}
  \left[
    1+
    2\sum_{k=1}^{n-1}
      \left(1-\frac{k}{n}\right)\rho_f(k)
  \right].
  \label{eq:stoch-markov-average-variance}
\end{align}
若$\sum_{k\geq1}|\rho_f(k)|<\infty$，括号中的量趋向
\[
  \tau_f
  =
  1+2\sum_{k=1}^{\infty}\rho_f(k),
\]
称为积分自相关时间。用$n_{\mathrm{eff}}\approx n/\tau_f$表达
\term{自相关有效样本量}（autocorrelation effective sample size），只是把该函数
均值的渐近方差换算成独立样本尺度。它依赖$f$，也不同于重要性权重的有效样本量。

对例~\ref{ex:stoch-two-state-chain}，式
\eqref{eq:stoch-two-state-autocorrelation}给出
\[
  \tau_Y=1+2\sum_{k=1}^{\infty}0.7^k
  =\frac{1+0.7}{1-0.7}=\frac{17}{3}.
\]
因而长链对均值$\mu$提供的信息量约为$3n/17$个独立观测，而不是$n$个。若自相关为负，
$\tau_f$也可能小于$1$；有效样本量是特定估计量的方差等价口径，不是链中“真正保留下来”
的样本条数。

\subsection{一般状态空间需要返回条件}

有限链可用矩阵、路径和周期描述可达性。连续或一般状态空间中，单点通常概率为零，“每个
状态互相到达”不再是合适表述。固定$(\mathcal{Z},\mathcal{A})$上的非零参考测度
$\varphi$。常用的$\varphi$不可约性要求：对每个满足$\varphi(B)>0$的
$B\in\mathcal{A}$以及每个初态$z\in\mathcal{Z}$，链从$z$出发都有正概率在有限时间内
到达$B$；若只研究某个不可约类，就先把$\mathcal{Z}$限制为该类。在此基础上，
\term{Harris常返性}（Harris recurrence）要求这些正测度集合以概率$1$被反复访问；
\term{正Harris常返性}（positive Harris recurrence）再要求存在不变概率分布。对
$\pi$可积的$f$，正Harris常返性提供一般状态空间遍历大数定律的标准入口。

这里“反复访问”带有概率$1$的量词：从$\mathcal{Z}$中的任意初态出发，每个满足
$\varphi(B)>0$的集合$B$都几乎必然被访问无穷多次。仅有“以正概率最终到达”还不够，
因为过程仍可能以正概率逃离后永不返回。正Harris常返链对其唯一不变概率$\pi$满足
式~\eqref{eq:stoch-markov-ergodic-average}的一般状态空间版本，但函数必须
$\pi$可积；若只知道局部可积或某个初态下样本值有限，不能推出该时间平均结论。

\term{几何遍历性}（geometric ergodicity）进一步要求到平稳分布的距离以几何速度下降，
例如存在$0<r<1$和依赖初态的有限$C(z)$，使
\[
  \lVert P^t(z,\cdot)-\pi\rVert_{\mathrm{TV}}
  \leq C(z)r^t.
\]
它比“存在平稳分布”或“时间平均收敛”更强，常与额外矩条件一起用于中心极限定理、Poisson
方程和随机逼近中的相关噪声控制。几何遍历是方便而有力的充分条件，不是所有遍历链都必须
满足的定义。

这些条件也解释了相关平均的边界。链若可约，长期平均可能由初始封闭类决定；链若只保持
$\pi$而不能充分移动，有限预算误差可能很大；观测函数若不可积，平稳分布存在也不能定义
所求均值。式~\eqref{eq:stoch-mean-root}在Markov观测下仍有意义，需要的是链的长期平均
确实识别$\mu$，而不是把相关的$Y_t$假装成独立样本。

\section{简单鞅、停时与序贯集中}
\label{sec:stoch-martingales-stopping-sequential-concentration}

\subsection{滤过与不可预测增量}

Markov链保留了当前状态对未来的预测作用。另一种更适合分析累计噪声的结构，是只要求每个
新增误差在给定过去后均值为零。在概率空间
$(\Omega,\mathcal{F},\mathbb{P})$上，令每个$\mathcal{F}_t$都是
$\mathcal{F}$的子$\sigma$代数，并且
\[
  \mathcal{F}_0\subseteq\mathcal{F}_1\subseteq\cdots\subseteq\mathcal{F}.
\]
这列递增的子$\sigma$代数称为\term{滤过}（filtration），可以理解为随时间增长的信息。
随机变量$X_t$若对
$\mathcal{F}_t$可测，就称它在时刻$t$已经可知；过程$(H_t)$若$H_t$对
$\mathcal{F}_{t-1}$可测，就称为\term{可预测过程}（predictable process）。这里的
“可预测”不是说$H_t$没有随机性，而是说它可以由第$t$次新噪声到来之前的信息确定。

可积过程$(M_t)$若适应于$(\mathcal{F}_t)$，并满足
\begin{equation}
  \mathbb{E}[M_t\mid\mathcal{F}_{t-1}]=M_{t-1},
  \label{eq:stoch-martingale-definition}
\end{equation}
就称为\term{鞅}（martingale）。增量$X_t=M_t-M_{t-1}$满足
\begin{equation}
  \mathbb{E}[X_t\mid\mathcal{F}_{t-1}]=0,
  \label{eq:stoch-martingale-difference}
\end{equation}
称为\term{鞅差}（martingale difference）。无条件零均值不够：若$X_t$可以由过去
预测，累计和仍可能形成系统漂移；式~\eqref{eq:stoch-martingale-difference}排除的正是
这种可利用的条件平均。

\begin{example}[公平随机游走与按过去下注]
  \label{ex:stoch-predictable-betting}
  设$\varepsilon_t$独立取$\{-1,1\}$，两种取值概率均为$1/2$，并令
  $\mathcal{F}_t=\sigma(\varepsilon_1,\ldots,\varepsilon_t)$。公平随机游走
  $S_t=\sum_{s=1}^{t}\varepsilon_s$是鞅。

  现在允许下注额$H_t$根据过去结果选择，只要求$H_t$对
  $\mathcal{F}_{t-1}$可测，并使$H_t\varepsilon_t$可积。财富过程
  \[
    M_t=M_0+\sum_{s=1}^{t}H_s\varepsilon_s
  \]
  仍是鞅，因为
  \[
    \mathbb{E}[H_t\varepsilon_t\mid\mathcal{F}_{t-1}]
    =
    H_t\mathbb{E}[\varepsilon_t\mid\mathcal{F}_{t-1}]
    =0.
  \]
  策略可以追涨、反向或在某些历史后停止下注，但只要它不能偷看$\varepsilon_t$，就不能
  改变下一笔收益的条件均值。若$H_t$依赖当前尚未揭晓的$\varepsilon_t$，上面的等式便
  失去依据。
\end{example}

对平方可积鞅差，定义\term{可预测二次变差}（predictable quadratic variation）
\begin{equation}
  V_t=\sum_{s=1}^{t}\mathbb{E}[X_s^2\mid\mathcal{F}_{s-1}].
  \label{eq:stoch-predictable-quadratic-variation}
\end{equation}
它在每一步使用新增量到来之前可知的条件方差，因而能够记录实际噪声规模。与之不同，
$\sum_{s=1}^{t}X_s^2$是观察后的实现二次变差。Freedman界使用前者；在其他序贯估计中
也可以围绕后者构造界，但不能不加说明地互换。

\subsection{停时与可选停止的条件}

随机时刻$\tau$若满足$\{\tau\leq t\}\in\mathcal{F}_t$对每个$t$成立，就称为
\term{停时}（stopping time）。到时刻$t$为止已有的信息足以判断是否已经停止；“最后一次
达到最大值的时刻”通常需要看到未来，不一定是停时。

\begin{theorem}[有界停时的可选停止]
  \label{thm:stoch-bounded-optional-stopping}
  设$(M_t,\mathcal{F}_t)$为可积鞅，$\tau$为满足$\tau\leq n$几乎必然成立的停时。
  则
  \[
    \mathbb{E}[M_\tau]=\mathbb{E}[M_0].
  \]
\end{theorem}

\begin{proof}
写$X_t=M_t-M_{t-1}$。由于$\tau$为停时，事件
$\{\tau\geq t\}=\{\tau\leq t-1\}^{\mathrm{c}}$属于$\mathcal{F}_{t-1}$，故其指示
函数是可预测的。又因为$\tau\leq n$，
\[
  M_\tau
  =
  M_0+\sum_{t=1}^{n}\mathbb{I}\{\tau\geq t\}X_t.
\]
对每一项使用塔式性质与鞅差条件，
\[
  \mathbb{E}\!\left[
    \mathbb{I}\{\tau\geq t\}X_t
  \right]
  =
  \mathbb{E}\!\left[
    \mathbb{I}\{\tau\geq t\}
    \mathbb{E}[X_t\mid\mathcal{F}_{t-1}]
  \right]
  =0.
\]
有限求和后即得结论。
\end{proof}

定理的证明清楚标出了有界性用途：它把随机长度的和改写为固定的有限和。无界停时需要
额外条件，例如鞅一致可积、停止后的变量由某个可积变量控制，或增量有界且
$\mathbb{E}\tau<\infty$。这些条件不能省略。对从$0$开始的公平随机游走，令
$\tau=\inf\{t:S_t=1\}$。该停时几乎必然有限，却有$S_\tau=1$，所以
$\mathbb{E}S_\tau=1\neq0$；这里$\mathbb{E}\tau=\infty$，不能套用有界停时定理。

\subsection{Azuma--Hoeffding界只使用统一增量尺度}

若只知道每个鞅差的范围，可以把独立变量的指数矩方法逐步条件化。

\begin{theorem}[Azuma--Hoeffding不等式]
  \label{thm:stoch-azuma-hoeffding}
  设$(X_t,\mathcal{F}_t)_{t=1}^{n}$为鞅差，并且几乎必然有
  $a_t\leq X_t\leq b_t$，其中$a_t,b_t$为确定常数。令
  $S_n=\sum_{t=1}^{n}X_t$。则对任意$x>0$，
  \begin{equation}
    \mathbb{P}(S_n\geq x)\leq\exp\left\{
      -\frac{2x^2}{\sum_{t=1}^{n}(b_t-a_t)^2}
    \right\}.
    \label{eq:stoch-azuma-hoeffding}
  \end{equation}
\end{theorem}

\begin{proof}
条件Hoeffding引理给出，对任意$\lambda>0$，
\[
  \mathbb{E}[
    \exp(\lambda X_t)
    \mid\mathcal{F}_{t-1}]
  \leq
  \exp\left\{
    \frac{\lambda^2(b_t-a_t)^2}{8}
  \right\}.
\]
从$t=n$开始反复使用塔式性质，得到
\[
  \mathbb{E}\exp(\lambda S_n)
  \leq
  \exp\left\{
    \frac{\lambda^2}{8}
    \sum_{t=1}^{n}(b_t-a_t)^2
  \right\}.
\]
Markov不等式于是给出
\[
  \mathbb{P}(S_n\geq x)
  \leq
  \exp\left\{
    -\lambda x+
    \frac{\lambda^2}{8}
    \sum_{t=1}^{n}(b_t-a_t)^2
  \right\}.
\]
对$\lambda$优化，取
$\lambda=4x/\sum_t(b_t-a_t)^2$，便得到式
\eqref{eq:stoch-azuma-hoeffding}。
\end{proof}

这个证明没有要求$X_t$彼此独立，只要求给定过去后均值为零，并且条件范围受控。代价是
它把每一步都按最坏范围计费。若多数时刻的条件方差很小，下一结果会给出更紧的尺度。

\subsection{Freedman界利用实际累积的条件方差}

\begin{theorem}[标量Freedman不等式]
  \label{thm:stoch-freedman}
  设$(X_t,\mathcal{F}_t)_{t\geq1}$为鞅差，且$|X_t|\leq b$几乎必然成立。
  令
  \[
    S_t=\sum_{s=1}^{t}X_s,\qquad
    V_t=\sum_{s=1}^{t}
      \mathbb{E}[X_s^2\mid\mathcal{F}_{s-1}].
  \]
  则对任意$x,v>0$，
  \begin{equation}
    \mathbb{P}\left(\exists t\geq1:S_t\geq x,\ V_t\leq v\right)
    \leq\exp\left\{
      -\frac{x^2}{2(v+bx/3)}
    \right\}.
    \label{eq:stoch-freedman}
  \end{equation}
\end{theorem}

\begin{proof}
对$0<\lambda<3/b$，标量不等式
\[
  \mathrm{e}^{\lambda u}
  \leq
  1+\lambda u+
  \frac{\lambda^2u^2}{2(1-\lambda b/3)},
  \qquad |u|\leq b,
  \]
与$\mathbb{E}[X_t\mid\mathcal{F}_{t-1}]=0$共同给出
\[
  \mathbb{E}[
    \mathrm{e}^{\lambda X_t}
    \mid\mathcal{F}_{t-1}]
  \leq
  \exp\left\{
    \psi(\lambda)
    \mathbb{E}[X_t^2\mid\mathcal{F}_{t-1}]
  \right\},
  \quad
  \psi(\lambda)=
  \frac{\lambda^2}{2(1-\lambda b/3)}.
\]
因此
\[
  L_t(\lambda)
  =
  \exp\{\lambda S_t-\psi(\lambda)V_t\},
  \qquad L_0(\lambda)=1,
\]
是非负超鞅。

令
\[
  \tau
  =
  \inf\{t\geq1:S_t\geq x,\ V_t\leq v\}.
  \]
对任意整数$n$，$\tau\wedge n$有界。对非负超鞅应用有界停时论证，
$\mathbb{E}L_{\tau\wedge n}(\lambda)\leq1$。在事件$\{\tau\leq n\}$上，
  \[
  L_\tau(\lambda)\geq\exp\{\lambda x-\psi(\lambda)v\}.
  \]
故
\[
  \mathbb{P}(\tau\leq n)\leq\exp\{-\lambda x+\psi(\lambda)v\}.
  \]
令$n\to\infty$，再取
$\lambda=x/(v+bx/3)$；该值位于$(0,3/b)$，代入后得到式
\eqref{eq:stoch-freedman}。
\end{proof}

若$X_1,\ldots,X_n$相互独立、中心化且$|X_t|\leq b$，并且
$\sum_{t=1}^{n}\operatorname{Var}(X_t)\leq v$，则自然滤过下的$X_t$构成鞅差，
且$V_n\leq v$确定成立。在式~\eqref{eq:stoch-freedman}中固定$t=n$便得到
\[
  \mathbb{P}(S_n\geq x)
  \leq
  \exp\left\{
    -\frac{x^2}{2(v+bx/3)}
  \right\}.
\]
这正是式~\eqref{eq:probability-bernstein-inequality}的同常数单侧Bernstein型界。
因此，Freedman界把独立变量的方差敏感界推广到了条件方差随历史变化的序贯情形。

Freedman界同时控制所有可能时刻，但事件中还限定了$V_t\leq v$。若希望得到随观测方差
自动变化的单一边界，可以把$V_t$的可能范围分成几何区间，对各区间使用不同$v$并分配
失败概率；更精细的方法直接混合不同$\lambda$。无论采用哪种构造，序贯有效性都来自
非负超鞅，而不是把固定时刻界中的$n$替换为随机停时。

这一点可由\term{Ville不等式}概括：若$L_t$为非负超鞅且
$\mathbb{E}L_0\leq1$，则
\begin{equation}
  \mathbb{P}\left(\sup_{t\geq0}L_t\geq\frac1\delta\right)
  \leq\delta.
  \label{eq:stoch-ville}
\end{equation}
证明只需令$\tau$为首次越过$1/\delta$的时刻，对$\tau\wedge n$使用有界停时，再令
$n\to\infty$。由此反演得到的\term{置信序列}（confidence sequence）在同一个高概率
事件上覆盖所有时刻，因而允许根据当前数据决定何时查看或停止。固定$n$的置信区间只承诺
预先指定时刻的覆盖；反复查看直到区间排除某个值，会改变错误概率。

同一个上穿论证还解释了后文为什么可以使用非负超鞅的极限。设$(L_t)$为非负超鞅且
$\mathbb{E}L_0<\infty$。对任意有理数$0<a<b$，记$U_n[a,b]$为
$L_0,\ldots,L_n$从不高于$a$到不低于$b$的完整上穿次数。在每次不高于$a$时持有一单位
过程、达到$b$时卖出；由于$(L_t)$是超鞅，这个可预测交易策略的期望收益不大于$0$。
比较每次完整交易至少$b-a$的路径收益与最后一次未完成交易，得到Doob上穿不等式
\[
  (b-a)U_n[a,b]
  \leq
  \text{交易收益}+(L_n-a)^-,
  \qquad x^-:=\max\{-x,0\}.
\]
其中负部$(L_n-a)^-$至多为$a$，而交易收益的期望不大于$0$。取期望得到
\[
  (b-a)\mathbb{E}U_n[a,b]
  \leq
  \mathbb{E}(L_n-a)^-
  \leq
  a.
\]
令$n\to\infty$可知每个正有理区间几乎必然只被上穿有限次。若$L_t$的下极限严格小于上
极限，就能在二者之间选到一个正有理区间并产生无穷多次上穿，形成矛盾。因此$L_t$几乎必然
存在扩展实数极限；再由Fatou引理与
$\sup_t\mathbb{E}L_t\leq\mathbb{E}L_0<\infty$排除极限为正无穷。于是非负超鞅几乎
必然收敛到有限随机变量。这个结论只保证路径极限存在；若还要交换极限与期望，仍需一致
可积等额外条件。

\section{自归一化集中与自适应设计}
\label{sec:stoch-self-normalized-adaptive-design}

\subsection{为什么误差要按设计矩阵衡量}

标量累计和把每次噪声视为同一个方向。在线性观测中，第$t$次查询带有特征
$\bm{x}_t\in\mathbb{R}^d$，观测为
\begin{equation}
  y_t=\langle\bm{x}_t,\bm{\theta}^{\star}\rangle+\eta_t.
  \label{eq:stoch-adaptive-linear-observation}
\end{equation}
允许$\bm{x}_t$根据$\mathcal{F}_{t-1}$中的历史结果选择，但要求它在$\eta_t$揭晓前
已经确定，也就是$\bm{x}_t$可预测；噪声$\eta_t$对$\mathcal{F}_t$可测。若噪声对
某个$R>0$满足条件次高斯界
\begin{equation}
  \mathbb{E}[
    \exp(\gamma\eta_t)
    \mid\mathcal{F}_{t-1}]
  \leq
  \exp\left(\frac{\gamma^2R^2}{2}\right),
  \qquad \gamma\in\mathbb{R},
  \label{eq:stoch-conditional-subgaussian}
\end{equation}
则它给定过去后均值为零，并且尾部由尺度$R$控制。

累计噪声和正则化设计矩阵分别记为
\begin{equation}
  \bm{S}_t=\sum_{s=1}^{t}\bm{x}_s\eta_s,
  \qquad
  \bm{V}_t=\bm{V}_0+\sum_{s=1}^{t}\bm{x}_s\bm{x}_s^{\top},
  \quad\bm{V}_0\succ\bm{0}.
  \label{eq:stoch-vector-sum-design}
\end{equation}
可预测性与条件中心化共同保证
\[
  \mathbb{E}[\bm{x}_t\eta_t\mid\mathcal{F}_{t-1}]
  =
  \bm{x}_t
  \mathbb{E}[\eta_t\mid\mathcal{F}_{t-1}]
  =\bm{0}.
\]
因此$(\bm{S}_t)$是向量鞅。这里没有要求$\bm{x}_t$为确定向量，也没有要求不同
$\bm{x}_t$彼此独立；关键是它在当前噪声出现前已经由过去确定。

某个方向被反复观测时，$\bm{V}_t$沿该方向增大；从未被询问的方向只剩初始正则化。
因此，自然误差尺度不是统一的欧氏球，而是
\[
  \lVert\bm{S}_t\rVert_{\bm{V}_t^{-1}}
  =\sqrt{\bm{S}_t^{\top}\bm{V}_t^{-1}\bm{S}_t}.
\]
设计提供的信息越多，逆矩阵在相应方向上的权重越小。这种由随机设计自身归一化误差的形式，
称为\term{自归一化}（self-normalized）；它与第~\ref{chap:probability-theory}章中用
随机权重和归一化分母构造的自归一化重要性估计不是同一个对象。

\subsection{混合指数超鞅给出向量界}

\begin{theorem}[向量自归一化集中界]
  \label{thm:stoch-vector-self-normalized}
  在式~\eqref{eq:stoch-adaptive-linear-observation}至
  \eqref{eq:stoch-vector-sum-design}的条件下，对任意$\delta\in(0,1)$，以至少
  $1-\delta$的概率，对所有$t\geq0$同时有
  \begin{equation}
    \lVert\bm{S}_t\rVert_{\bm{V}_t^{-1}}
    \leq
    R\sqrt{
      2\log\left(
        \frac{\det(\bm{V}_t)^{1/2}}
             {\det(\bm{V}_0)^{1/2}\delta}
      \right)
    }.
    \label{eq:stoch-vector-self-normalized-bound}
  \end{equation}
\end{theorem}

\begin{proof}
固定任意$\bm{u}\in\mathbb{R}^d$。由于$\bm{x}_t$对
$\mathcal{F}_{t-1}$可测，在式~\eqref{eq:stoch-conditional-subgaussian}中取
$\gamma=\langle\bm{u},\bm{x}_t\rangle$，得到
\[
  \mathbb{E}\!\left[
    \exp\left\{
      \langle\bm{u},\bm{x}_t\rangle\eta_t
      -\frac{R^2}{2}
       \langle\bm{u},\bm{x}_t\rangle^2
    \right\}
    \middle|\mathcal{F}_{t-1}
  \right]
  \leq1.
  \]
逐步相乘可知
\[
  L_t(\bm{u})
  =
  \exp\left\{
    \langle\bm{u},\bm{S}_t\rangle
    -\frac{R^2}{2}
     \lVert\bm{u}\rVert_{\bm{V}_t-\bm{V}_0}^2
  \right\}
\]
是非负超鞅。

一个固定$\bm{u}$只能控制一个投影方向。为同时处理所有方向，取独立于数据的随机向量
$\bm{U}\sim\mathcal{N}(\bm{0},R^{-2}\bm{V}_0^{-1})$，并对
$L_t(\bm{U})$关于$\bm{U}$积分。非负性允许交换条件期望与积分，所以混合过程
$\overline L_t=\mathbb{E}_{\bm{U}}L_t(\bm{U})$仍是非负超鞅，且
$\overline L_0=1$。把高斯密度与指数中的二次项合并，配方得到
\begin{align*}
  \overline L_t
  &=
  \left(
    \frac{\det(\bm{V}_0)}
         {\det(\bm{V}_t)}
  \right)^{1/2}
  \exp\left\{
    \frac{1}{2R^2}
    \lVert\bm{S}_t\rVert_{\bm{V}_t^{-1}}^2
  \right\}.
\end{align*}
这里行列式比来自两个高斯二次型的归一化常数，而指数项来自配方后的中心移动。

对$\overline L_t$使用式~\eqref{eq:stoch-ville}，以至少$1-\delta$的概率，对所有$t$
都有$\overline L_t<1/\delta$。取对数并整理，便得到式
\eqref{eq:stoch-vector-self-normalized-bound}。
\end{proof}

证明中的四个条件各有明确用途。条件次高斯性构造每一步指数超鞅；特征可预测性允许在条件
期望中把$\bm{x}_t$视为已知；$\bm{V}_0\succ\bm{0}$使高斯混合与逆矩阵从$t=0$起有
定义；行列式比累计各方向相对于初始尺度增长了多少。行列式不是把每个访问次数简单相加，
因为同一方向的重复观测与新方向的首次观测提供的信息不同。

\subsection{从噪声界到置信椭球}

取$\bm{V}_0=\lambda\bm{I}$，$\lambda>0$，并定义岭回归估计
\begin{equation}
  \widehat{\bm{\theta}}_t
  =\bm{V}_t^{-1}\sum_{s=1}^{t}\bm{x}_s y_s.
  \label{eq:stoch-ridge-estimator}
\end{equation}
由式~\eqref{eq:stoch-adaptive-linear-observation}，
\[
  \bm{V}_t
  (\widehat{\bm{\theta}}_t-\bm{\theta}^{\star})
  =
  \bm{S}_t-\lambda\bm{\theta}^{\star}.
  \]
若还知道$\lVert\bm{\theta}^{\star}\rVert_2\leq B$，三角不等式给出
\begin{align}
  \lVert
    \widehat{\bm{\theta}}_t-\bm{\theta}^{\star}
  \rVert_{\bm{V}_t}
  &\leq
  \lVert\bm{S}_t\rVert_{\bm{V}_t^{-1}}
  +
  \lambda
  \lVert\bm{\theta}^{\star}\rVert_{\bm{V}_t^{-1}}
  \notag\\
  &\leq
  R\sqrt{
    2\log\left(
      \frac{\det(\bm{V}_t)^{1/2}}
           {\lambda^{d/2}\delta}
    \right)
  }
  +\sqrt{\lambda}B
  =:\beta_t(\delta).
  \label{eq:stoch-adaptive-confidence-ellipsoid}
\end{align}
第二个不等式使用$\bm{V}_t\succeq\lambda\bm{I}$。因此
\[
  \mathcal{C}_t(\delta)
  =
  \left\{
    \bm{\theta}:
    \lVert\bm{\theta}-\widehat{\bm{\theta}}_t
    \rVert_{\bm{V}_t}
    \leq\beta_t(\delta)
  \right\}
\]
在同一个概率至少为$1-\delta$的事件上，对所有$t$都包含真实参数。这是置信椭球而不是
欧氏球：观测充分的方向较窄，信息不足的方向较宽。

\begin{example}[二维自适应查询]
  \label{ex:stoch-two-dimensional-adaptive-design}
  设$d=2$、$\lambda=1$，每次只能在
  $\bm{e}_1=(1,0)^\top$与$\bm{e}_2=(0,1)^\top$中选择一个观测方向。算法可以在第
  $t$步比较当前椭球在两个方向上的宽度，并选择较宽的方向；这个选择只依赖过去，因而
  $\bm{x}_t$可预测。

  若前四次选择中$\bm{e}_1$出现三次、$\bm{e}_2$出现一次，则
  \[
    \bm{V}_4
    =
    \begin{pmatrix}
      4&0\\
      0&2
    \end{pmatrix},
    \qquad
    \det(\bm{V}_4)=8.
  \]
  对任意方向$\bm{x}$，椭球给出的预测误差满足
  \[
    |\langle
      \bm{x},
      \widehat{\bm{\theta}}_4-\bm{\theta}^{\star}
    \rangle|
    \leq
    \beta_4(\delta)
    \lVert\bm{x}\rVert_{\bm{V}_4^{-1}}.
  \]
  因而两个坐标方向的半宽分别与$1/2$和$1/\sqrt2$成比例。第一个方向访问更多，区间
  更窄；下一次转向第二个方向是由已有信息决定的自适应设计。定理仍然适用，因为它在每次
  新噪声产生之前固定当前方向，并在一个事件上同时控制全部时刻。
\end{example}

这个例子也说明为什么不能把独立样本界中的$n$换成某个动作的访问次数。参数误差是向量，
不同特征方向可能相关；置信宽度由$\bm{x}^{\top}\bm{V}_t^{-1}\bm{x}$决定，而不是只由
某个坐标被访问多少次决定。行列式项又要为算法可能探索的全部方向支付复杂度。

\subsection{自适应有效不等于任意数据依赖都有效}

自归一化界允许$\bm{x}_t$依赖$\mathcal{F}_{t-1}$，不允许它依赖当前$\eta_t$后再伪装成
预先选择。若传感器先看到当前噪声再决定是否记录，$\bm{x}_t\eta_t$的条件平均可能不再为
零，指数超鞅的第一步就会失败。类似地，模型失配使
$\mathbb{E}[y_t\mid\mathcal{F}_{t-1},\bm{x}_t]$不等于
$\langle\bm{x}_t,\bm{\theta}^{\star}\rangle$时，椭球会同时包含系统偏差与随机误差，
不能继续解释为围绕同一个真实参数的置信集合。

条件次高斯性也不是有限方差的同义词。重尾噪声可能需要截断、稳健均值或只使用有限矩的
其他序贯工具；直接套用式~\eqref{eq:stoch-vector-self-normalized-bound}会低估尾部。
$\lambda$保证初期矩阵可逆并编码参数范数先验，但过大的$\lambda$会增加
$\sqrt{\lambda}B$偏差项。置信椭球控制统计误差，不保证选择策略会访问所有对任务重要的
方向；探索是否充分仍要由设计规则与任务目标另行证明。

\section{随机逼近与Markov噪声}
\label{sec:stoch-approximation-markov-noise}

\subsection{从保存平均到持续求根}

若全部观测都可保存，均值求根可以在$t$步后直接使用样本平均。流式算法则希望每到一个
观测只更新当前估计。对式~\eqref{eq:stoch-mean-root}，Robbins--Monro型递推为
\begin{equation}
  \theta^{(t+1)}
  =\theta^{(t)}+\alpha_t\bigl(Y_{t+1}-\theta^{(t)}\bigr),
  \qquad t\geq0,
  \label{eq:stoch-robbins-monro-mean}
\end{equation}
其中确定数$\alpha_t>0$为步长。新观测高于当前估计时参数上移，低于当前估计时参数下移；
若
$\mathbb{E}[Y_{t+1}\mid\mathcal{F}_t]=\mu$，则条件平均更新方向为
$\mu-\theta^{(t)}$，正好指向根。Markov观测一般不满足这个条件中心化假设，后文将用
Poisson方程分离相关波动
\scite{math-robbins1951stochastic}

若$\alpha_t=1/(t+1)$，并从$\theta^{(0)}$任意值开始，则第一次更新完全采用$Y_1$，
归纳可得
\[
  \theta^{(t)}
  =
  \frac1t\sum_{s=1}^{t}Y_s,
  \qquad t\geq1.
\]
一般步长不再给出等权平均。较大的近期步长让算法更快响应新观测，也让噪声保留得更久；
递减步长试图在持续移动与逐渐消除随机波动之间取得平衡。

更一般的随机逼近写成
\begin{equation}
  \bm{\theta}^{(t+1)}
  =\bm{\theta}^{(t)}+\alpha_t\left[
    \overline{\bm{h}}(\bm{\theta}^{(t)})
    +\bm{\xi}_{t+1}
    +\bm{r}_{t+1}
  \right].
  \label{eq:stoch-general-stochastic-approximation}
\end{equation}
$\overline{\bm{h}}$是平均漂移，$\bm{\xi}_{t+1}$是给定过去后均值为零的噪声，
$\bm{r}_{t+1}$是近似、截断或分布尚未平衡产生的偏差。算法要收敛到
$\overline{\bm{h}}(\bm{\theta})=\bm{0}$的稳定根，必须分别控制这三部分，不能只验证
单步估计“看起来无偏”。

\subsection{步长条件分别控制移动与噪声}

经典步长条件为
\begin{equation}
  \sum_{t=0}^{\infty}\alpha_t=\infty,
  \qquad
  \sum_{t=0}^{\infty}\alpha_t^2<\infty.
  \label{eq:stoch-robbins-monro-stepsizes}
\end{equation}
第一项保证算法保留无限的累计移动能力；若步长和有限，初始误差可能只被缩小有限倍。
第二项使有限方差噪声的加权方差可求和。取
$\alpha_t=c/(t+t_0)^\gamma$时，通常
$1/2<\gamma\leq1$满足这两个条件。固定步长不满足平方可和条件，适合追踪变化目标或维持
稳态误差，却一般不保证迭代精确收敛到一个点。

下面先给出独立或鞅差噪声下的完整代表性证明。证明只用一个特殊的超鞅收敛引理。

\begin{lemma}[带可求和误差的非负超鞅]
  \label{lem:stoch-almost-supermartingale}
  设$(Z_t)$与$(A_t)$都是关于$(\mathcal{F}_t)$的非负可积适应过程，
  $(b_t)$为非负确定数列，且以下不等式几乎必然成立：
  \[
    \mathbb{E}[Z_{t+1}\mid\mathcal{F}_t]
    \leq Z_t-A_t+b_t.
  \]
  若$\sum_tb_t<\infty$，则$Z_t$几乎必然收敛到有限随机变量，且
  $\sum_tA_t<\infty$几乎必然成立。
\end{lemma}

\begin{proof}
令$r_t=\sum_{s=t}^{\infty}b_s$，并取$U_t=Z_t+r_t$。由于
$r_t=b_t+r_{t+1}$，
\[
  \mathbb{E}[U_{t+1}\mid\mathcal{F}_t]
  \leq U_t-A_t.
  \]
$(U_t)$是非负超鞅，故几乎必然收敛；对上式取期望并从$0$到$n$求和，得到
$\mathbb{E}\sum_{t=0}^{n}A_t\leq\mathbb{E}U_0$。令$n\to\infty$，由单调收敛定理，
$\mathbb{E}\sum_{t=0}^{\infty}A_t\leq\mathbb{E}U_0<\infty$，所以
$\sum_tA_t<\infty$几乎必然成立。又因$r_t\to0$，$Z_t=U_t-r_t$也收敛。
\end{proof}

\begin{theorem}[均值求根的几乎必然收敛]
  \label{thm:stoch-mean-root-convergence}
  设$\mathcal{F}_t$包含$\theta^{(0)},Y_1,\ldots,Y_t$，且
  $\mathbb{E}[(\theta^{(0)}-\mu)^2]<\infty$。假设
  \[
    \mathbb{E}[Y_{t+1}\mid\mathcal{F}_t]=\mu,\qquad
    \mathbb{E}[(Y_{t+1}-\mu)^2\mid\mathcal{F}_t]\leq\sigma^2,
  \]
  步长满足式~\eqref{eq:stoch-robbins-monro-stepsizes}，并且从某时刻起
  $0<\alpha_t\leq1$。则递推~\eqref{eq:stoch-robbins-monro-mean}满足
  $\theta^{(t)}\to\mu$几乎必然成立。
\end{theorem}

\begin{proof}
记$e_t=\theta^{(t)}-\mu$，
$\xi_{t+1}=Y_{t+1}-\mu$。递推变为
\[
  e_{t+1}
  =
  (1-\alpha_t)e_t+\alpha_t\xi_{t+1}.
  \]
取$T$使得对所有$t\geq T$都有$0<\alpha_t\leq1$。由初值二阶矩有限、条件二阶矩
有界以及上式递推归纳可知，每个有限时刻的$\mathbb{E}e_t^2$都有限，特别是
$\mathbb{E}e_T^2<\infty$。以下只考虑$t\geq T$；舍去有限前缀不改变极限，也不改变
两组步长级数的收敛或发散性质。

条件于$\mathcal{F}_t$展开平方，交叉项因
$\mathbb{E}[\xi_{t+1}\mid\mathcal{F}_t]=0$消失：
\begin{align*}
  \mathbb{E}[e_{t+1}^2\mid\mathcal{F}_t]
  &=
  (1-\alpha_t)^2e_t^2
  +
  \alpha_t^2
  \mathbb{E}[\xi_{t+1}^2\mid\mathcal{F}_t]\\
  &\leq
  e_t^2-\alpha_t e_t^2+\alpha_t^2\sigma^2,
\end{align*}
其中使用$0<\alpha_t\leq1$得到
$(1-\alpha_t)^2\leq1-\alpha_t$。

令$\mathcal{G}_k=\mathcal{F}_{T+k}$。对关于$(\mathcal{G}_k)$的移位序列，在
引理~\ref{lem:stoch-almost-supermartingale}中取
$Z_k=e_{T+k}^2$、$A_k=\alpha_{T+k}e_{T+k}^2$、
$b_k=\alpha_{T+k}^2\sigma^2$。上述二阶矩归纳还验证了$Z_k$与$A_k$的可积性。平方可和
条件给出$\sum_{k=0}^{\infty}b_k<\infty$，所以$e_t^2$几乎必然收敛，并且
$\sum_{t=T}^{\infty}\alpha_te_t^2<\infty$。若其极限在某条样本路径上为$L>0$，则
充分大的$t$都有
$e_t^2\geq L/2$，从而
\[
  \sum_{t=T}^{\infty}\alpha_te_t^2
  \geq
  \frac{L}{2}\sum_{t=T}^{\infty}\alpha_t
  =
  \infty,
  \]
产生矛盾。因此$L=0$，即$\theta^{(t)}\to\mu$。
\end{proof}

这个证明把两个步长条件的职责分开了。平方可和使噪声项
$\alpha_t^2\sigma^2$可累计控制；步长和发散则排除非零极限误差。若噪声具有无限条件
方差，或条件均值带有不可消失的偏差，上述递推不等式不再成立。若迭代是向量且平均漂移
非线性，还要另行保证迭代不会逃向无穷远。

\subsection{累计步长把随机更新连接到平均ODE}

定义算法时间
\begin{equation}
  \tau_0=0,\qquad \tau_t=\sum_{s=0}^{t-1}\alpha_s.
  \label{eq:stoch-algorithmic-time}
\end{equation}
在区间$[\tau_t,\tau_{t+1}]$上对$\bm{\theta}^{(t)}$与
$\bm{\theta}^{(t+1)}$作线性插值，就得到连续曲线。若任意固定算法时间窗内的累计噪声与
累计偏差趋于零，且$\overline{\bm{h}}$足够规则，那么这条插值曲线在晚期近似平均常微分
方程
\begin{equation}
  \dot{\bm{\theta}}(\tau)=\overline{\bm{h}}(\bm{\theta}(\tau)).
  \label{eq:stoch-mean-ode}
\end{equation}
这就是随机逼近的\term{ODE方法}（ODE method）：不要求单步随机更新接近零，而要求在
与步长匹配的时间尺度上，随机波动不会改变平均轨迹的极限行为
\scite{math-borkar2008stochastic}

对均值求根，
\[
  \dot{\theta}=\mu-\theta.
\]
从初值$\theta(0)$出发的解为
\begin{equation}
  \theta(\tau)=\mu+
  \bigl(\theta(0)-\mu\bigr)\mathrm{e}^{-\tau}.
  \label{eq:stoch-mean-root-ode-solution}
\end{equation}
因此$\mu$是全局渐近稳定平衡。等价地，取
$W(\theta)=(\theta-\mu)^2/2$，沿ODE轨迹有
\[
  \frac{\mathrm{d}}{\mathrm{d}\tau}W(\theta(\tau))
  =-(\theta(\tau)-\mu)^2\leq0,
\]
且只有在根处等于零。这解释了定理
\ref{thm:stoch-mean-root-convergence}中误差平方为何是自然的势函数。

上面的标量例子只需显式解出ODE，便能看出平均漂移把轨迹拉向\(\mu\)。一般向量递推还
需要更完整的条件：平均漂移的正则性保证ODE解适定，迭代有界性把分析限制在可控区域，
加权噪声与累计偏差则要在固定算法时间窗内消失。满足这些条件后，才能把离散随机插值与
平均ODE的长期行为联系起来；平均ODE若有多个稳定对象，也不能据此断言算法到达全局
最优解。有关流、稳定性与Lyapunov方法的系统定义和判据见第
\ref{chap:dynamical-systems-and-state-estimation}章，本章只保留这个可直接求解的
标量ODE作为前导。

迭代有界性不能从局部稳定直接推出。常见处理包括证明一个径向无界Lyapunov函数，也就是
当$\lVert\bm{\theta}\rVert_2\to\infty$时函数值趋于无穷，并验证它在远处具有负漂移；
另一种做法是把每一步投影到紧凸集。投影只保证迭代有界，还会把平均动力学改成带边界
约束的投影ODE。只有当这个投影ODE存在相应的稳定约束平衡，并且噪声、步长等条件同时
成立时，算法才可能收敛到该平衡。原方程的根位于集合外这一事实本身，既不能保证收敛，
也不能把可能的极限直接解释成原方程根在边界上的投影。

\subsection{Poisson方程分离Markov相关波动}

定理~\ref{thm:stoch-mean-root-convergence}假设
$Y_{t+1}-\mu$给定$\mathcal{F}_t$后均值为零。对Markov链，这通常不成立：
例~\ref{ex:stoch-two-state-chain}中，
$\mathbb{E}[Y_{t+1}-\mu\mid Z_t]=0.7(Y_t-\mu)$。不过，相关噪声仍可以借助
\term{Poisson方程}（Poisson equation）拆成鞅差与近似望远镜和。

先固定参数，设链的转移核为$P$、平稳分布为$\pi$，并令
$H(\theta,z)$为一次随机更新。平均漂移是
\begin{equation}
  \overline h(\theta)=\int H(\theta,z)\,\pi(\mathrm{d}z).
  \label{eq:stoch-markov-mean-field}
\end{equation}
对每个$\theta$，若存在函数$g_\theta$满足
\begin{equation}
  g_\theta(z)-Pg_\theta(z)
  =
  H(\theta,z)-\overline h(\theta),
  \qquad
  Pg_\theta(z)
  =
  \int g_\theta(z')P(z,\mathrm{d}z'),
  \label{eq:stoch-poisson-equation}
\end{equation}
则对$Z_{t+1}\sim P(Z_t,\cdot)$，
\begin{align}
  H(\theta,Z_t)-\overline h(\theta)
  &=
  g_\theta(Z_t)-g_\theta(Z_{t+1})
  +
  \Delta M_{t+1}(\theta),
  \label{eq:stoch-poisson-decomposition}\\
  \Delta M_{t+1}(\theta)
  &:=
  g_\theta(Z_{t+1})-Pg_\theta(Z_t).
  \notag
\end{align}
最后一项给定$\mathcal{F}_t$后均值为零，是鞅差。前两项在固定$\theta$、固定步长时会
望远镜消去；步长和参数缓慢变化时，分部求和会留下由
$|\alpha_t-\alpha_{t-1}|$以及
$\lVert\theta^{(t)}-\theta^{(t-1)}\rVert$控制的余项。

\begin{lemma}[固定有限链的加权Poisson分解]
  \label{lem:stoch-finite-chain-poisson-sum}
  设$P$是有限状态空间上的不可约非周期核，且不随参数变化。令参数始终位于紧集$K$，
  $H(\theta,z)$在$K\times\mathcal{Z}$上有界，并关于$\theta$一致Lipschitz。考虑递推
  \[
    \theta^{(t+1)}
    =
    \theta^{(t)}+\alpha_tH(\theta^{(t)},Z_t),
  \]
  其中正步长单调递减并满足式~\eqref{eq:stoch-robbins-monro-stepsizes}。则中心化
  Markov噪声的加权和
  \[
    \sum_{t=0}^{n}\alpha_t
    \bigl[
      H(\theta^{(t)},Z_t)-\overline h(\theta^{(t)})
    \bigr]
  \]
  几乎必然收敛；特别地，它在任意向后移动的有限算法时间窗上的尾和趋于$0$。
\end{lemma}

\begin{proof}
有限不可约非周期链以几何速度趋于$\pi$。对
$q_\theta(z)=H(\theta,z)-\overline h(\theta)$，有
$\pi q_\theta=0$，故级数
\[
  g_\theta(z)=\sum_{k=0}^{\infty}P^kq_\theta(z)
\]
绝对且一致收敛，并满足
$g_\theta-Pg_\theta=q_\theta$。$H$的一致有界性与Lipschitz性质还给出常数
$G,L_g<\infty$，使
\[
  |g_\theta(z)|\leq G,\qquad
  |g_\theta(z)-g_{\theta'}(z)|
  \leq L_g|\theta-\theta'|.
\]
这一步是有限链和固定核提供的关键统一控制。

把式~\eqref{eq:stoch-poisson-decomposition}代入从$m$到$n$的加权和。鞅差部分为
\[
  \sum_{t=m}^{n}\alpha_t
  \Delta M_{t+1}(\theta^{(t)}).
\]
它的增量有界，条件二阶矩之和由常数倍$\sum_t\alpha_t^2$控制，所以该鞅在
$L^2$中有界并几乎必然收敛。剩余部分可作离散分部求和：
\begin{align*}
  &\sum_{t=m}^{n}\alpha_t
  \left[
    g_{\theta^{(t)}}(Z_t)-g_{\theta^{(t)}}(Z_{t+1})
  \right]\\
  &\quad=
  \alpha_mg_{\theta^{(m)}}(Z_m)
  -\alpha_ng_{\theta^{(n)}}(Z_{n+1})\\
  &\qquad+
  \sum_{t=m+1}^{n}
  \left[
    \alpha_tg_{\theta^{(t)}}(Z_t)
    -\alpha_{t-1}g_{\theta^{(t-1)}}(Z_t)
  \right].
\end{align*}
中间每项的绝对值不超过
\[
  G(\alpha_{t-1}-\alpha_t)
  +
  \alpha_{t-1}L_g
  |\theta^{(t)}-\theta^{(t-1)}|.
\]
递推与$H$有界给出
$|\theta^{(t)}-\theta^{(t-1)}|\leq C\alpha_{t-1}$。第一部分按单调步长望远镜求和，
第二部分由$\sum_t\alpha_t^2<\infty$控制，两个端点又因$\alpha_t\to0$而消失。
因此望远镜余项和鞅项都具有有限极限，所述加权和收敛。收敛级数的尾和趋于$0$，自然也
覆盖任意固定算法时间长度对应的有限索引窗。
\end{proof}

例~\ref{ex:stoch-two-state-chain}可以显式求解这个方程。取
$H(\theta,z)=z-\theta$，则
$\overline h(\theta)=\mu-\theta$，方程右侧为$z-\mu$。由式
\eqref{eq:stoch-two-state-eigenfunction}，
$P(z-\mu)=0.7(z-\mu)$，故可取
\[
  g_\theta(z)=\frac{z-\mu}{1-0.7}=\frac{10}{3}(z-\mu).
\]
这里$g_\theta$恰好与$\theta$无关且有界，相关噪声的分解尤其简单。这个计算说明
Poisson方程不是把Markov观测改造成独立观测，而是把累计相关误差重新组织为可由鞅工具
控制的部分和边界余项。

在一般状态空间中，Poisson方程解的存在、矩界与参数正则性需要由链的遍历性质保证。
若转移核还随参数变化，即
\[
  Z_{t+1}\sim P_{\bm{\theta}^{(t)}}(Z_t,\cdot),
\]
就必须在迭代访问的参数区域上统一控制$P_{\bm{\theta}}$的遍历速度、平稳分布
$\pi_{\bm{\theta}}$以及$g_{\bm{\theta}}$随参数的变化。只证明每个固定参数下的链分别
遍历，不足以控制一个持续改变转移核的算法。常用的充分条件包括紧参数区域上的统一几何
遍历、更新函数与Poisson解的适当矩界和Lipschitz连续性，再配合
式~\eqref{eq:stoch-robbins-monro-stepsizes}与迭代有界性。

在这些条件下，Markov噪声随机逼近的证明路线与独立噪声相接：Poisson分解产生加权鞅和，
平方可和步长控制其波动；望远镜项及参数变化余项在有限算法时间窗上消失；插值过程因而
跟随
\[
  \dot{\bm{\theta}}
  =\overline{\bm{h}}(\bm{\theta})
  =\int\bm{H}(\bm{\theta},z)\,\pi_{\bm{\theta}}(\mathrm{d}z).
\]
最终收敛到哪里仍由这个平均ODE的稳定结构决定。混合只帮助把相关观测替换为平均漂移加
可控噪声，不会把一个不稳定平衡变成稳定平衡。

\subsection{双时间尺度与仍需核对的边界}

有些算法同时更新快变量$\bm{u}^{(t)}$与慢变量$\bm{v}^{(t)}$：
\begin{align*}
  \bm{u}^{(t+1)}
  &=\bm{u}^{(t)}+\beta_t\bigl[
    \bm{g}(\bm{u}^{(t)},\bm{v}^{(t)})
    +\text{噪声}
  \bigr],\\
  \bm{v}^{(t+1)}
  &=\bm{v}^{(t)}+\alpha_t\bigl[
    \bm{f}(\bm{u}^{(t)},\bm{v}^{(t)})
    +\text{噪声}
  \bigr].
\end{align*}
若两组步长都满足Robbins--Monro条件，且
$\alpha_t/\beta_t\to0$，则$\bm{u}$使用相对较大的步长，是快变量；在它的时间尺度上，
$\bm{v}$近似冻结。证明先要求固定$\bm{v}$时快ODE具有稳定平衡
$\bm{u}^{\star}(\bm{v})$，再把慢变量的平均ODE写成
\[
  \dot{\bm{v}}
  =
  \bm{f}(\bm{u}^{\star}(\bm{v}),\bm{v}).
\]
步长比只建立时间尺度分离，不保证任一ODE稳定，也不代替噪声、Markov混合和联合有界性
条件。若两组噪声是鞅差，通常还要求其条件二阶矩在迭代访问区域内统一有界；若噪声来自
Markov链，则要分别为快、慢更新验证相应的Poisson分解与统一遍历控制。平衡映射
$\bm{u}^{\star}(\bm{v})$还需在慢变量访问区域内足够正则，才能把快过程的跟踪误差传入
慢平均漂移。

\begin{example}[分解同一个未知均值的快慢更新]
  \label{ex:stoch-two-timescale-mean-splitting}
  设观测满足
  $\mathbb{E}[Y_{t+1}\mid\mathcal{F}_t]=\mu$，考虑两个标量递推
  \begin{align*}
    u^{(t+1)}
    &=
    u^{(t)}
    +\beta_t\bigl(Y_{t+1}-u^{(t)}-v^{(t)}\bigr),\\
    v^{(t+1)}
    &=
    v^{(t)}
    +\alpha_t\bigl(u^{(t)}-v^{(t)}\bigr),
  \end{align*}
  其中两组步长分别满足Robbins--Monro条件，且
  $\alpha_t/\beta_t\to0$。固定慢变量$v$时，快平均ODE为
  $\dot u=\mu-v-u$，其唯一全局稳定平衡是
  $u^\star(v)=\mu-v$。把这个平衡代入慢更新，得到
  \[
    \dot v=u^\star(v)-v=\mu-2v,
  \]
  所以慢平衡为$v^\star=\mu/2$，相应
  $u^\star(v^\star)=\mu/2$。这个例子中，快变量先估计当前$v$尚未解释的均值部分，
  慢变量再调整两部分的分配；若
  $\mathbb{E}[(Y_{t+1}-\mu)^2\mid\mathcal{F}_t]$统一有界、初值二阶矩有限且联合
  迭代几乎必然有界，时间尺度论证指向
  $(u^{(t)},v^{(t)})\to(\mu/2,\mu/2)$。

  若只知道$\alpha_t/\beta_t\to0$，却没有快平衡的稳定性、映射
  $u^\star(v)$的正则性、慢ODE的稳定性或迭代有界性，上述结论均不能推出。若两组步长
  同阶，这个线性例子仍可另作联合系统分析，但不能再用“慢变量近似冻结”的证明。
\end{example}

随机逼近还有几类常见边界。若$\sum_t\alpha_t<\infty$，算法可能在到达根之前停止有效
移动；若$\sum_t\alpha_t^2=\infty$且噪声持续存在，随机波动可能无法消失。固定步长通常
围绕稳定根形成稳态分布，误差随步长缩小，但目标变化时反而可能比递减步长更能跟踪。
观测具有重尾、系统偏差或缓慢变化的均值时，需要稳健化、偏差控制或跟踪分析。参数平均
可以改善某些正则模型的渐近方差，却不能修复不稳定平均漂移或缺失的访问条件。

后续动力系统与状态估计章将在状态动力学已经建立后，进一步分析一般向量ODE、非对称
线性漂移与时序差分更新的稳定性。本章提供的接口是：先写出平均漂移，再证明噪声在算法
时间上可忽略，最后核对平均动力学的稳定结构。任何一步缺失，都不能仅凭损失曲线平滑或
样本数量增加声称随机更新已经收敛。

\section{本章小结}
\label{sec:stochastic-processes-and-approximation-summary}

当观测沿时间相关或由过去自适应选择时，独立样本分析失效的位置并不相同。Markov链用
转移核描述当前状态怎样影响下一步；平稳分布只保证从该分布出发后保持不变，不可约性、
周期与返回条件决定能否从指定初态接近它。时间平均还受到自相关影响，所以记录数不能直接
解释为独立信息量。有限链的路径条件也不能原样搬到一般状态空间；正Harris常返和几何遍历
分别提供长期平均与更强误差分析的入口。

鞅把关注点从状态依赖转向条件不可预测性。可预测策略可以根据过去改变每一步的尺度，只要
新增量给定过去后均值为零，累计过程仍可用指数超鞅分析。Azuma--Hoeffding按最坏增量范围
计费，Freedman进一步利用可预测条件方差。停时与Ville不等式解释了序贯保证的来源，也
说明固定时刻概率界为何不能直接支持反复查看后再停止。

带可预测特征的向量鞅需要按设计矩阵归一化。高斯混合把所有投影方向合成一个非负超鞅，
由此得到同时覆盖全部时刻的自归一化界和置信椭球。椭球会随已经选择的方向收缩，因此能够
适应过去驱动的设计；它仍要求当前特征在当前噪声揭晓前确定，并要求条件尾部与线性模型
假设成立。

均值求根把这些误差工具放入持续更新。步长和发散保留到达根所需的累计移动，平方和收敛
控制有限方差噪声，平均漂移的稳定性决定更新指向哪里，迭代有界性保证局部分析不会失去
适用范围。Markov噪声可借Poisson方程分成鞅差与近似望远镜项，再由算法时间插值接入平均
ODE；受参数控制的链和双时间尺度还需统一遍历、正则性与稳定条件。由此，相关平均、
自适应置信和随机逼近共同回答了章首问题：数据可以依赖过去，但证明必须准确保留条件
平均、信息增长和长期动力学之间的关系。第~\ref{chap:mathematical-statistics}章将
回到有限观测的推断问题，区分可识别性、估计误差、区间与检验，并说明抽样制度怎样决定
这些结论所针对的总体。
